\section{Experiments}
\label{sec:experiments}

\subsection{Evaluation Protocol}

We evaluate 39 RGB-D loop-closure cases from 18 Kimera-Multi
sequences~\cite{tian23arxiv_kimeramultiexperiments} and 12 LiDAR cases from
six KITTI sequences~\cite{Geiger2013IJRR}. Every method reconstructs the same
temporal interval and uses the same GT: a 5 m tube around the official
Kimera-Multi trajectory or GT-registered KITTI points within the mapper's 25 m
ray range. Predictions outside the GT frame are aligned by a RANSAC--Umeyama
$\mathrm{Sim}(3)$ fitted only to timestamp-matched trajectory positions; no
mesh-to-GT registration is used. Inputs are replayed offline, but map updates
remain sequential: each corrected TSDF replaces the active Voxblox volume and
fusion continues to the next loop closure.

\noindent\textbf{Baselines.}
Uncorrected TSDF uses the pre-correction trajectory. Open3D~\cite{zhou2018open3dmodernlibrary3d}
and Wavemap~\cite{reijgwart2023wavemap} replay corrected poses into an
independent TSDF and a wavelet occupancy map, respectively; Voxgraph
~\cite{reijgwart2020voxgraph} optimizes rigid TSDF submaps. Voxblox
reintegration~\cite{Oleynikova_2017} is the representation-matched reference.
WarpVox instead uses the uncorrected TSDF, sensor-frame support, and the
pre/post-correction trajectories. For deterministic offline evaluation, Stage
1 is reproduced by replaying only recorded keyframes after the pre-mesh is
available; online, the same associations accumulate as keyframes arrive.

We compute precision and recall by bidirectional nearest-neighbour tests and
report F@10/25/50 at 0.10/0.25/0.50 m. All methods share each case's crop,
spatial region, two-million-point sample, random seed, and thresholds. Main
TSDF experiments use 0.25 m voxels and 1.0 m truncation; Wavemap uses a matched
0.25 m finest level. Core times are measured on an AMD Ryzen 7 9700X with 64 GB
RAM from in-memory inputs; bag reading, disk I/O, serialization, visualization,
evaluation, and evaluation-only Marching Cubes are excluded.

\subsection{Reconstruction Fidelity and Update Cost}

Table~\ref{tab:broad_main} gives the uniform 0.25 m comparison. WarpVox is
the closest method to Voxblox reintegration at every threshold, with absolute
deviations below 0.27 pp on Kimera-Multi and 0.39 pp on KITTI, and has the best
F@25 and F@50 among non-reference methods on both datasets.

\begin{table}[t]
\centering
\caption{\textbf{Reconstruction accuracy and mean core update time at
0.25 m.} F-scores are percentages; lower time is better.}
\label{tab:broad_main}
\scriptsize
\setlength{\tabcolsep}{2.1pt}
\begin{tabular}{llrrrr}
\toprule
Data & Method & F@10 & F@25 & F@50 & Core (s) \\
\midrule
\multirow{6}{*}{K-Multi (39)}
 & Uncorrected TSDF & 7.00 & 19.29 & 34.91 & -- \\
 & Open3D & 8.13 & 22.94 & 41.30 & -- \\
 & Voxgraph & 8.27 & 22.78 & 39.49 & 5.54 \\
 & Wavemap & 1.73 & 11.54 & 30.91 & 8.32 \\
 & Voxblox reintegration & 10.93 & 28.71 & \textbf{48.67} & 109.25 \\
 & \textbf{WarpVox} & \textbf{11.13} & \textbf{28.77} & 48.40 & 2.43 \\
\midrule
\multirow{6}{*}{KITTI (12)}
 & Uncorrected TSDF & 7.13 & 27.91 & 49.22 & -- \\
 & Open3D & \textbf{13.50} & 38.58 & 61.02 & -- \\
 & Voxgraph & 9.38 & 31.89 & 52.49 & 4.53 \\
 & Wavemap & 5.23 & 33.10 & 62.55 & 2.68 \\
 & Voxblox reintegration & 13.21 & \textbf{44.29} & \textbf{67.45} & 15.73 \\
 & \textbf{WarpVox} & 13.00 & 43.98 & 67.06 & 4.89 \\
\bottomrule
\end{tabular}
\vspace{1pt}

\parbox{\columnwidth}{\scriptsize Times for unchanged baselines are retained
from their canonical runs; WarpVox is remeasured with the final implementation.}
\end{table}

WarpVox is $45.1\times$ faster than reintegration on Kimera-Multi and
$3.22\times$ on KITTI. The smaller KITTI gain follows from cheap sparse-LiDAR
replay versus WarpVox's larger active surface and output volume, while its TSDF
remains the closest alternative to the reference.

Table~\ref{tab:runtime_breakdown} separates the four stages. For
Kimera-Multi, the unchanged Stage 1--2 measurements are combined with the
remeasured final Stage 3--4 implementation. For KITTI, all four stages are
remeasured; its Stage 3 mean comprises 2.457 s of ARAP deformation and
0.148 s of topology processing.

\begin{table}[t]
\centering
\caption{\textbf{Mean WarpVox core-time breakdown in seconds.}}
\label{tab:runtime_breakdown}
\scriptsize
\setlength{\tabcolsep}{2.3pt}
\begin{tabular}{lrrrrr}
\toprule
Setting & Stage 1 & Stage 2 & Stage 3 & Stage 4 & Total \\
\midrule
K-Multi, 25 cm (39) & 0.460 & 0.400 & 1.375 & 0.190 & 2.425 \\
KITTI, 25 cm (12)   & 0.310 & 0.069 & 2.605 & 1.907 & 4.890 \\
K-Multi, 10 cm (11) & 4.138 & 0.250 & 5.052 & 0.993 & 10.432 \\
\bottomrule
\end{tabular}
\vspace{1pt}

\parbox{\columnwidth}{\scriptsize Stage 1: association; Stage 2: target
aggregation; Stage 3: deformation and topology processing; Stage 4: TSDF
reconstruction. Component means are rounded independently.}
\end{table}

As a resolution-transfer check, the same 11 Kimera-Multi cases reconstructed
at 0.10 m obtain F@10/25/50 of 17.13/39.76/57.83 with WarpVox, compared with
17.21/39.68/57.12 for Voxblox reintegration. WarpVox requires 10.432 s of
core computation versus 192.548 s for reintegration, an $18.5\times$
reduction. The close quantitative agreement is complemented by the
qualitative comparisons in Fig.~\ref{fig:qualitative_ral}.

\begin{figure*}[t]
\centering
\setlength{\tabcolsep}{0pt}
\renewcommand{\arraystretch}{0.2}
\newcommand{\qa}[1]{\includegraphics[width=.245\linewidth,trim=150 150 150 150,clip]{#1}}
\newcommand{\qb}[1]{\includegraphics[width=.245\linewidth,trim=250 70 250 70,clip]{#1}}
\newcommand{\qc}[1]{\includegraphics[width=.245\linewidth,trim=200 20 200 20,clip]{#1}}
\newcommand{\qd}[1]{\includegraphics[width=.245\linewidth,trim=100 20 100 20,clip]{#1}}
\begin{tabular}{cccc}
\qa{Figures/qualitative/5_voxblox_pre.png} &
\qa{Figures/qualitative/5_voxblox_post.png} &
\qa{Figures/qualitative/5_voxgraph_post_sep.png} &
\qa{Figures/qualitative/5_warped_tsdf.png} \\[2pt]
\qb{Figures/qualitative/28_voxblox_pre.png} &
\qb{Figures/qualitative/28_voxblox_post.png} &
\qb{Figures/qualitative/28_voxgraph_post_sep.png} &
\qb{Figures/qualitative/28_warped_tsdf.png} \\[2pt]
\qc{Figures/qualitative/40_voxblox_pre.png} &
\qc{Figures/qualitative/40_voxblox_post.png} &
\qc{Figures/qualitative/40_voxgraph_post_sep.png} &
\qc{Figures/qualitative/40_warped_tsdf.png} \\[2pt]
\qd{Figures/qualitative/35_voxblox_pre.png} &
\qd{Figures/qualitative/35_voxblox_post.png} &
\qd{Figures/qualitative/35_voxgraph_post_sep.png} &
\qd{Figures/qualitative/35_warped_tsdf.png} \\[4pt]
\qc{Figures/qualitative/76_voxblox_pre.png} &
\qc{Figures/qualitative/76_voxblox_post.png} &
\qc{Figures/qualitative/76_voxgraph_post_sep.png} &
\qc{Figures/qualitative/76_warped_tsdf.png} \\[2pt]
\qd{Figures/qualitative/77_voxblox_pre.png} &
\qd{Figures/qualitative/77_voxblox_post.png} &
\qd{Figures/qualitative/77_voxgraph_post_sep.png} &
\qd{Figures/qualitative/77_warped_tsdf.png} \\[2pt]
\small Voxblox$_{\mathrm{Pre}}$ &
\small Voxblox reintegration &
\small Voxgraph &
\small \textbf{WarpVox}
\end{tabular}
\caption{\textbf{Qualitative reconstruction comparison at 0.25 m.} Four
Kimera-Multi loop-closure cases are shown in the top four rows and two KITTI
cases in the bottom two rows; each row uses the same view across methods.}
\label{fig:qualitative_ral}
\end{figure*}

\subsection{Event-Triggered Comparison with Kimera-PGMO}

Kimera-PGMO continuously maintains a MeshFrontend deformation graph, whereas
WarpVox is triggered by a global trajectory revision. We compare final states
on ten complete Kimera-Multi bags using the same 5 cm Voxblox pre-mesh and
final VIO-to-PGMO correction. PGMO additionally uses timestamped graph
associations (3 m spacing, three neighbours, 20 s horizon); WarpVox constructs
support from recorded keyframes and VIO poses. When asynchronous replays vary,
we select the complete trajectory--mesh pair with lowest ATE, independently of
mesh accuracy, and apply its correction to both methods. Because the native
pre-mesh has no source TSDF, WarpVox uses mesh-only Stage 4; this comparison
therefore tests surface correction, not source-support transfer.

\begin{table}[t]
\centering
\caption{\textbf{Mean final-state comparison over ten complete Kimera-Multi
bags at 5 cm.} F-scores are percentages.}
\label{tab:pgmo_full10}
\scriptsize
\setlength{\tabcolsep}{2.2pt}
\begin{tabular}{lrrrrr}
\toprule
Method & F@5 & F@10 & F@25 & F@50 & Core (s) \\
\midrule
Kimera-PGMO mesh & 1.96 & 6.42 & 17.48 & 31.54 & 435.97 \\
\textbf{WarpVox (mesh-only TSDF)} & \textbf{1.97} & \textbf{6.53} &
\textbf{18.72} & \textbf{34.57} & \textbf{47.02} \\
\bottomrule
\end{tabular}
\vspace{1pt}

\parbox{\columnwidth}{\scriptsize Kimera-PGMO time accumulates graph
optimization and mesh deformation over all callbacks; WarpVox time covers one
Ray, Ceres, ARAP, and mesh-only TSDF update.}
\end{table}

One WarpVox update improves mean F@25/F@50 by 1.24/3.03 pp while using
47.02 s versus PGMO's 435.97 s accumulated over all callbacks. It has higher
F@50 on eight of ten bags ($-0.44$ to $+11.74$ pp). Thus WarpVox offers a
lower-cost final-state correction, while PGMO provides intermediate updates and
frontend temporal associations.

\subsection{Ablation Study}

Table~\ref{tab:ablation} evaluates every variant on the same ten selected
0.25 m Kimera-Multi checkpoints and separates deformation scale from the
Stage 3 and Stage 4 design changes.

\begin{table}[t]
\centering
\caption{\textbf{Ablation of the deformation and TSDF reconstruction stages
at 0.25 m.}}
\label{tab:ablation}
\scriptsize
\setlength{\tabcolsep}{1.1pt}
\begin{tabular}{@{}llrrrrr@{}}
\toprule
Component & Variant & $n$ & F@10 & F@25 & F@50 & Core (s) \\
\midrule
\multirow{2}{*}{Scale}
 & Full-mesh ARAP & 10 & 10.278 & 26.749 & 45.490 & 3.323 \\
 & 60K QEM-proxy ARAP$^\dagger$ & 10 & 10.277 & 26.717 & 45.436 & 1.408 \\
\midrule
\multirow{2}{*}{Stage 3}
 & Grid-proxy ARAP & 10 & 10.254 & 26.464 & 44.877 & 1.471 \\
 & Robust QEM-proxy ARAP$^\dagger$ & 10 & 10.277 & 26.717 & 45.436 & 1.408 \\
\midrule
\multirow{2}{*}{Stage 4}
 & Mesh-only TSDF & 10 & 9.982 & 26.400 & 45.068 & 0.098 \\
 & Source-supported TSDF & 10 & 10.155 & 26.616 & 45.392 & 0.160 \\
\bottomrule
\end{tabular}
\vspace{1pt}

\parbox{\columnwidth}{\scriptsize $^\dagger$Identical final Stage-3
configuration shown in both pairwise comparisons. Core time measures only the
ablated component.}
\end{table}

The 60K QEM proxy is $2.36\times$ faster than full-mesh ARAP with only
$-0.032/-0.054$ pp at F@25/F@50. Against the grid proxy, robust QEM gains
$0.023/0.253/0.559$ pp and is slightly faster; all mesh-identity checks pass.
Source-supported reconstruction adds 0.062 s, gains $0.173/0.216/0.324$ pp,
and suppresses unsupported surfaces at open boundaries.

\subsection{Extension to Gaussian Primitives}

To test representation transfer, we apply the same targets to a
SplaTAM~\cite{keetha2024splatam} map while changing only Gaussian centres.
On one sequence, F@25/F@50 rises from 17.23/35.25 to 23.12/40.59 and depth L1
falls from 173.93 to 142.54 cm in about 21 s; a corrected-pose rebuild obtains
23.25/41.01 and 141.37 cm in about 2 h. PSNR/MS-SSIM/LPIPS improve from
11.96/0.512/0.430 to 12.59/0.526/0.411. This isolates the target mechanism:
Gaussian splats~\cite{kerbl2023gaussians} still lack the signed-distance and
observed-free-space semantics that keep TSDF/ESDF mapping as WarpVox's primary
robotics setting.